\documentclass[a4paper,11pt]{article}
\usepackage[utf8]{inputenc}
\usepackage[T1]{fontenc}
\usepackage{amsmath,amssymb}
\usepackage[a4paper,left=2cm,right=2cm,top=3.6cm,bottom=2.3cm,headheight=1.1cm,headsep=0.6cm,footskip=1cm]{geometry}
\usepackage{xcolor}
\usepackage{tikz}
\usepackage{fancyhdr}
\usepackage{lastpage}
\usepackage{tabularx}
\usepackage{array}
\usepackage[most]{tcolorbox}
\usepackage{enumitem}
\usepackage{needspace}

\definecolor{navy}{HTML}{1F3A5F}
\definecolor{teal}{HTML}{0F8B7D}
\definecolor{softbg}{HTML}{F4F7FA}
\newcommand{\dis}{\displaystyle}

\setlength{\parindent}{0pt}
\renewcommand{\arraystretch}{1.35}

% ---------- header berwarna navy ----------
\pagestyle{fancy}
\fancyhf{}
\renewcommand{\headrulewidth}{0pt}
\fancyhead[L]{%
  \begin{tikzpicture}[remember picture,overlay]
    \fill[navy] ([yshift=-1.2cm]current page.north west) rectangle ([yshift=-3.15cm]current page.north east);
    \fill[teal] ([yshift=-3.15cm]current page.north west) rectangle ([yshift=-3.25cm]current page.north east);
  \end{tikzpicture}%
  \color{white}\bfseries\large Latihan Soal Barisan dan Deret}
\fancyhead[R]{\color{white}\small Diunduh dari website \textbf{Math1729}\ \,|\,\ 4 Oktober 2026}
\fancyfoot[L]{\small\color{navy}Math1729 --- Barisan dan Deret}
\fancyfoot[R]{\small\color{navy}Halaman \thepage\ dari \pageref{LastPage}}

% ---------- komponen ----------
\newcounter{no}
\newcommand{\bagian}[2]{\par\Needspace{6\baselineskip}\vspace{1.0em}\noindent
  \begin{tcolorbox}[colback=navy,colframe=navy,coltext=white,boxrule=0pt,arc=2pt,left=6pt,right=6pt,top=3pt,bottom=3pt]
  \textbf{#1}\hfill{\small #2}\end{tcolorbox}\setcounter{no}{0}\par\vspace{-0.3em}}
\newcommand{\tingkat}[2]{\par\Needspace{7\baselineskip}\vspace{0.6em}\noindent
  {\color{teal}\rule{3pt}{1.1em}}\ {\color{navy}\bfseries #1}\ {\small\color{gray}--- #2}\par\vspace{0.1em}}
\newcommand{\soal}[2]{\par\vspace{0.75em}\stepcounter{no}\noindent
  \begin{minipage}[t]{1.8em}\textbf{\theno.}\end{minipage}%
  \begin{minipage}[t]{\dimexpr\linewidth-1.8em\relax}#1\par\vspace{0.35em}#2\end{minipage}\par}
\newcommand{\poin}[1]{\hfill{\small\color{navy}[#1 poin]}}


\begin{document}
\thispagestyle{fancy}

\begin{center}
  {\color{navy}\fontsize{22}{26}\selectfont\bfseries Latihan Soal Barisan dan Deret\par}
\end{center}
\vspace{2pt}
\begin{tcolorbox}[colback=softbg,colframe=navy!40,boxrule=0.6pt,arc=3pt,left=8pt,right=8pt,top=5pt,bottom=5pt]
\noindent\textbf{Nama} : \rule{4.6cm}{0.4pt}\hfill\textbf{Kelas} : \rule{2.2cm}{0.4pt}\hfill\textbf{Skor} : \rule{2.2cm}{0.4pt}
\end{tcolorbox}
\vspace{2pt}
\begin{tcolorbox}[colback=white,colframe=teal,boxrule=0.8pt,arc=3pt,left=8pt,right=8pt,top=4pt,bottom=4pt,title={\bfseries Petunjuk Pengerjaan},coltitle=white,fonttitle=\small]
\small
\begin{enumerate}[leftmargin=1.4em,itemsep=1pt,topsep=1pt]
\item Soal terdiri atas \textbf{20 pilihan ganda} (5 mudah, 5 sedang, 5 sulit, 5 HOTS), \textbf{5 isian singkat}, dan \textbf{5 uraian (essay)}. Alokasi waktu yang disarankan: 120 menit.
\item Skor: pilihan ganda 2 poin, isian singkat 4 poin, uraian 8 poin per soal (total 100 poin). Tidak ada pengurangan skor untuk jawaban salah.
\item Notasi: $a$ suku pertama, $b$ beda, $r$ rasio, $U_n$ suku ke-$n$, dan $S_n$ jumlah $n$ suku pertama.
\item Sebelum memilih rumus, tentukan dahulu pola yang tetap: \textbf{selisih} (aritmetika) atau \textbf{perbandingan} (geometri). Pada deret tak hingga, periksa syarat $|r|<1$.
\end{enumerate}
\end{tcolorbox}

\bagian{A. Pilihan Ganda}{20 soal $\times$ 2 poin = 40 poin. Pilih satu jawaban yang paling tepat.}
\tingkat{Tingkat Mudah}{soal 1--5, konsep dasar}
\soal{Suku ke-15 dari barisan aritmetika $100,93,86,79,\ldots$ adalah \dots.}{\begin{tabularx}{\linewidth}{@{}*{5}{X}@{}}\textbf{A.}\;$-5$ & \textbf{B.}\;$2$ & \textbf{C.}\;$9$ & \textbf{D.}\;$16$ & \textbf{E.}\;$23$\end{tabularx}}
\soal{Suku ke-7 dari barisan geometri $2,-6,18,-54,\ldots$ adalah \dots.}{\begin{tabularx}{\linewidth}{@{}*{5}{X}@{}}\textbf{A.}\;$-1.458$ & \textbf{B.}\;$-486$ & \textbf{C.}\;$729$ & \textbf{D.}\;$1.458$ & \textbf{E.}\;$2.916$\end{tabularx}}
\soal{Jumlah $20$ suku pertama deret $4+9+14+19+\cdots$ adalah \dots.}{\begin{tabularx}{\linewidth}{@{}*{5}{X}@{}}\textbf{A.}\;$1.030$ & \textbf{B.}\;$1.080$ & \textbf{C.}\;$1.130$ & \textbf{D.}\;$1.180$ & \textbf{E.}\;$1.260$\end{tabularx}}
\soal{Di antara $5$ dan $77$ disisipkan $11$ bilangan sehingga membentuk barisan aritmetika. Bilangan kelima yang disisipkan adalah \dots.}{\begin{tabularx}{\linewidth}{@{}*{5}{X}@{}}\textbf{A.}\;$23$ & \textbf{B.}\;$29$ & \textbf{C.}\;$35$ & \textbf{D.}\;$41$ & \textbf{E.}\;$47$\end{tabularx}}
\soal{Jumlah deret geometri tak hingga $18+6+2+\dis\frac23+\cdots$ adalah \dots.}{\begin{tabularx}{\linewidth}{@{}*{5}{X}@{}}\textbf{A.}\;$9$ & \textbf{B.}\;$12$ & \textbf{C.}\;$18$ & \textbf{D.}\;$24$ & \textbf{E.}\;$27$\end{tabularx}}
\tingkat{Tingkat Sedang}{soal 6--10, dua langkah atau lebih}
\soal{Pada barisan aritmetika diketahui $U_4=17$ dan $U_9=37$. Jumlah $12$ suku pertama barisan tersebut adalah \dots.}{\begin{tabularx}{\linewidth}{@{}*{5}{X}@{}}\textbf{A.}\;$324$ & \textbf{B.}\;$336$ & \textbf{C.}\;$348$ & \textbf{D.}\;$360$ & \textbf{E.}\;$372$\end{tabularx}}
\soal{Barisan geometri positif memenuhi $U_2=12$ dan $U_5=96$. Jumlah $7$ suku pertamanya adalah \dots.}{\begin{tabularx}{\linewidth}{@{}*{5}{X}@{}}\textbf{A.}\;$372$ & \textbf{B.}\;$378$ & \textbf{C.}\;$381$ & \textbf{D.}\;$384$ & \textbf{E.}\;$762$\end{tabularx}}
\soal{Gaji pertama seorang karyawan Rp4.000.000,00 per bulan dan naik Rp250.000,00 setiap bulan. Total gaji yang ia terima selama $24$ bulan pertama adalah \dots.}{\begin{tabularx}{\linewidth}{@{}*{3}{X}@{}}\textbf{A.}\;Rp132.000.000,00 & \textbf{B.}\;Rp150.000.000,00 & \textbf{C.}\;Rp165.000.000,00\\[3pt] \textbf{D.}\;Rp168.000.000,00 & \textbf{E.}\;Rp180.000.000,00 & \end{tabularx}}
\soal{Suku pertama suatu deret geometri tak hingga adalah $12$ dan jumlahnya $16$. Suku ke-$4$ deret tersebut adalah \dots.}{\begin{tabularx}{\linewidth}{@{}*{5}{X}@{}}\textbf{A.}\;$\dis\frac{3}{128}$ & \textbf{B.}\;$\dis\frac{3}{64}$ & \textbf{C.}\;$\dis\frac{3}{32}$ & \textbf{D.}\;$\dis\frac{3}{16}$ & \textbf{E.}\;$\dis\frac{3}{8}$\end{tabularx}}
\soal{Jumlah $n$ suku pertama suatu deret dirumuskan $S_n=2n^{2}+5n$. Suku ke-$12$ deret tersebut adalah \dots.}{\begin{tabularx}{\linewidth}{@{}*{5}{X}@{}}\textbf{A.}\;$43$ & \textbf{B.}\;$45$ & \textbf{C.}\;$47$ & \textbf{D.}\;$49$ & \textbf{E.}\;$51$\end{tabularx}}
\tingkat{Tingkat Sulit}{soal 11--15, bergaya UTBK-SNBT}
\soal{Jumlah $5$ suku pertama suatu deret aritmetika adalah $55$, dan jumlah $5$ suku berikutnya (suku ke-$6$ sampai ke-$10$) adalah $130$. Jumlah suku ke-$11$ sampai suku ke-$15$ adalah \dots.}{\begin{tabularx}{\linewidth}{@{}*{5}{X}@{}}\textbf{A.}\;$175$ & \textbf{B.}\;$190$ & \textbf{C.}\;$205$ & \textbf{D.}\;$215$ & \textbf{E.}\;$230$\end{tabularx}}
\soal{Diketahui barisan aritmetika dengan $U_2=-5$ dan $U_6=11$. Perhatikan pernyataan berikut.\par\smallskip\hspace*{0.3em}(1) Suku positif pertama adalah suku ke-$4$.\par\hspace*{0.3em}(2) $S_{10}=96$.\par\hspace*{0.3em}(3) Bilangan $2.027$ merupakan salah satu suku barisan.\par\hspace*{0.3em}(4) Terdapat dua suku barisan yang selisihnya $30$.\par\smallskip Pernyataan yang benar adalah \dots.}{\begin{tabularx}{\linewidth}{@{}X@{}}\textbf{A.}\;(1), (2), dan (3)\\[2pt] \textbf{B.}\;(1) dan (3)\\[2pt] \textbf{C.}\;(2) dan (4)\\[2pt] \textbf{D.}\;hanya (4)\\[2pt] \textbf{E.}\;semua pernyataan benar\end{tabularx}}
\soal{Rata-rata $10$ suku pertama suatu deret aritmetika adalah $20$, sedangkan rata-rata $20$ suku pertamanya adalah $30$. Jumlah $30$ suku pertama deret tersebut adalah \dots.}{\begin{tabularx}{\linewidth}{@{}*{5}{X}@{}}\textbf{A.}\;$1.080$ & \textbf{B.}\;$1.140$ & \textbf{C.}\;$1.200$ & \textbf{D.}\;$1.260$ & \textbf{E.}\;$1.320$\end{tabularx}}
\soal{Jumlah seluruh suku suatu deret geometri tak hingga positif adalah $36$, sedangkan jumlah suku-suku pada urutan ganjil ($U_1+U_3+U_5+\cdots$) adalah $27$. Nilai $U_3$ adalah \dots.}{\begin{tabularx}{\linewidth}{@{}*{5}{X}@{}}\textbf{A.}\;$2$ & \textbf{B.}\;$\dis\frac83$ & \textbf{C.}\;$4$ & \textbf{D.}\;$\dis\frac{16}{3}$ & \textbf{E.}\;$8$\end{tabularx}}
\soal{Suku ke-$2$, ke-$5$, dan ke-$14$ suatu barisan aritmetika membentuk barisan geometri. Jika suku pertama barisan aritmetika itu $4$, maka jumlah $10$ suku pertamanya adalah \dots.}{\begin{tabularx}{\linewidth}{@{}*{5}{X}@{}}\textbf{A.}\;$340$ & \textbf{B.}\;$360$ & \textbf{C.}\;$380$ & \textbf{D.}\;$400$ & \textbf{E.}\;$420$\end{tabularx}}
\tingkat{HOTS}{soal 16--20, penalaran dan konteks tidak rutin}
\soal{Jeruk ditumpuk berbentuk limas segitiga: lapisan teratas $1$ buah, lapisan kedua $3$ buah, lapisan ketiga $6$ buah, lapisan keempat $10$ buah, dan seterusnya (tiap lapisan berbentuk segitiga susun). Jika tumpukan terdiri atas $15$ lapisan, banyak jeruk seluruhnya adalah \dots.}{\begin{tabularx}{\linewidth}{@{}*{5}{X}@{}}\textbf{A.}\;$560$ & \textbf{B.}\;$600$ & \textbf{C.}\;$680$ & \textbf{D.}\;$720$ & \textbf{E.}\;$816$\end{tabularx}}
\soal{Diberikan barisan $3,5,9,15,23,\ldots$ Suku ke-$20$ barisan tersebut adalah \dots.}{\begin{tabularx}{\linewidth}{@{}*{5}{X}@{}}\textbf{A.}\;$341$ & \textbf{B.}\;$361$ & \textbf{C.}\;$380$ & \textbf{D.}\;$381$ & \textbf{E.}\;$383$\end{tabularx}}
\soal{Suatu deret aritmetika memiliki banyak suku ganjil. Jumlah suku-suku pada urutan ganjil adalah $56$ dan jumlah suku-suku pada urutan genap adalah $48$. Banyak suku deret tersebut adalah \dots.}{\begin{tabularx}{\linewidth}{@{}*{5}{X}@{}}\textbf{A.}\;$7$ & \textbf{B.}\;$9$ & \textbf{C.}\;$11$ & \textbf{D.}\;$13$ & \textbf{E.}\;$15$\end{tabularx}}
\soal{Nilai dari $\dis\sum_{k=1}^{99}\dis\frac{1}{\sqrt{k}+\sqrt{k+1}}$ adalah \dots.}{\begin{tabularx}{\linewidth}{@{}*{5}{X}@{}}\textbf{A.}\;$9$ & \textbf{B.}\;$10$ & \textbf{C.}\;$11$ & \textbf{D.}\;$50$ & \textbf{E.}\;$99$\end{tabularx}}
\soal{Sebuah persegi bersisi $8$ cm. Titik-titik tengah keempat sisinya dihubungkan sehingga terbentuk persegi kedua; dengan cara yang sama pada persegi kedua terbentuk persegi ketiga, dan seterusnya tanpa akhir. Jumlah keliling semua persegi tersebut adalah \dots\ cm.}{\begin{tabularx}{\linewidth}{@{}*{3}{X}@{}}\textbf{A.}\;$64$ & \textbf{B.}\;$32+32\sqrt2$ & \textbf{C.}\;$64+16\sqrt2$\\[3pt] \textbf{D.}\;$64+32\sqrt2$ & \textbf{E.}\;$128$ & \end{tabularx}}
\bagian{B. Isian Singkat}{5 soal $\times$ 4 poin = 20 poin. Tuliskan hanya jawaban akhir.}
\soal{Di antara $3$ dan $192$ disisipkan $5$ bilangan sehingga terbentuk barisan geometri dengan rasio positif. Jumlah \textbf{seluruh} suku barisan itu (termasuk $3$ dan $192$) adalah \dots}{\hfill Jawaban: \rule{4.5cm}{0.4pt}}
\soal{Nilai dari $\dis\sum_{k=5}^{20}(3k-2)$ adalah \dots}{\hfill Jawaban: \rule{4.5cm}{0.4pt}}
\soal{Populasi bakteri bertambah menjadi tiga kali lipat setiap $2$ jam. Jika populasi awal $200$ bakteri, banyak bakteri setelah $8$ jam adalah \dots}{\hfill Jawaban: \rule{4.5cm}{0.4pt}}
\soal{Jumlah semua bilangan asli di antara $100$ dan $500$ yang habis dibagi $7$ adalah \dots}{\hfill Jawaban: \rule{4.5cm}{0.4pt}}
\soal{Bilangan desimal berulang $1{,}272727\ldots$ jika dinyatakan sebagai pecahan paling sederhana adalah \dots}{\hfill Jawaban: \rule{4.5cm}{0.4pt}}
\bagian{C. Uraian (Essay)}{5 soal $\times$ 8 poin = 40 poin. Tuliskan langkah penyelesaian dengan lengkap.}
\soal{Diketahui deret aritmetika dengan $U_4=14$ dan $S_8=124$.\par\smallskip\hspace*{0.3em}(a) Tentukan suku pertama dan beda deret tersebut.\par\hspace*{0.3em}(b) Suku ke berapakah yang bernilai $83$?\par\hspace*{0.3em}(c) Tentukan bilangan asli $n$ terkecil sehingga $S_n>1.000$.}{\poin{8}}
\soal{Pada suatu deret geometri dengan suku-suku positif, jumlah $2$ suku pertama adalah $12$ dan jumlah $4$ suku pertama adalah $120$.\par\smallskip\hspace*{0.3em}(a) Tentukan suku pertama dan rasio deret tersebut.\par\hspace*{0.3em}(b) Tentukan rumus $U_n$ dan nilai $U_6$.\par\hspace*{0.3em}(c) Tentukan bilangan asli $n$ terkecil sehingga $S_n>1.000$.}{\poin{8}}
\soal{Sebuah bola dijatuhkan dari ketinggian $10$ m. Setiap kali memantul, bola mencapai $\dis\frac35$ dari tinggi pantulan sebelumnya.\par\smallskip\hspace*{0.3em}(a) Tentukan tinggi pantulan ke-$3$.\par\hspace*{0.3em}(b) Tentukan panjang seluruh lintasan bola sampai berhenti.\par\hspace*{0.3em}(c) Pada pantulan ke berapa tinggi bola pertama kali kurang dari $1$ m? Tunjukkan perhitungannya.}{\poin{8}}
\soal{Jumlah $n$ suku pertama suatu deret adalah $S_n=3n^{2}-n$.\par\smallskip\hspace*{0.3em}(a) Tentukan rumus suku ke-$n$ ($U_n$).\par\hspace*{0.3em}(b) Buktikan bahwa barisan tersebut aritmetika, lalu tentukan suku pertama dan bedanya.\par\hspace*{0.3em}(c) Hitung $U_{11}+U_{12}+\cdots+U_{20}$.}{\poin{8}}
\soal{Tiga bilangan positif membentuk barisan geometri naik dengan hasil kali $64$. Jika suku kedua ditambah $1$, ketiga bilangan itu membentuk barisan aritmetika.\par\smallskip\hspace*{0.3em}(a) Tentukan suku tengah barisan geometri tersebut.\par\hspace*{0.3em}(b) Tentukan rasio barisan geometri tersebut.\par\hspace*{0.3em}(c) Tentukan jumlah $10$ suku pertama barisan aritmetika yang terbentuk (dilanjutkan dengan pola yang sama).}{\poin{8}}
\newpage
\bagian{Kunci Jawaban dan Pembahasan Singkat}{Untuk guru dan pembahasan mandiri}
\par\vspace{0.6em}{\bfseries\color{navy}Kunci Pilihan Ganda}\par\vspace{0.3em}
\begin{tabularx}{\linewidth}{@{}*{10}{>{\centering\arraybackslash}X}@{}}\hline \textbf{1} & \textbf{2} & \textbf{3} & \textbf{4} & \textbf{5} & \textbf{6} & \textbf{7} & \textbf{8} & \textbf{9} & \textbf{10}\\ \hline B & D & A & C & E & A & E & C & D & E\\ \hline\end{tabularx}\par\vspace{0.4em}\begin{tabularx}{\linewidth}{@{}*{10}{>{\centering\arraybackslash}X}@{}}\hline \textbf{11} & \textbf{12} & \textbf{13} & \textbf{14} & \textbf{15} & \textbf{16} & \textbf{17} & \textbf{18} & \textbf{19} & \textbf{20}\\ \hline C & B & C & B & D & C & E & D & A & D\\ \hline\end{tabularx}
\par\vspace{0.8em}{\bfseries\color{navy}Pembahasan Pilihan Ganda}\par\vspace{0.2em}
\small\begin{enumerate}[leftmargin=2em,itemsep=2pt,topsep=2pt]
\item[\textbf{1.}] \textbf{(B)} $a=100$, $b=-7$. $U_{15}=100+14(-7)=100-98=2$.
\item[\textbf{2.}] \textbf{(D)} $a=2$, $r=\dis\frac{-6}{2}=-3$. $U_{7}=2(-3)^{6}=2\cdot729=1.458$ (pangkat genap, hasil positif).
\item[\textbf{3.}] \textbf{(A)} $a=4$, $b=5$. $S_{20}=\dis\frac{20}{2}\big(2(4)+19(5)\big)=10(103)=1.030$.
\item[\textbf{4.}] \textbf{(C)} Ada $11+1=12$ interval, sehingga $b=\dis\frac{77-5}{12}=6$. Bilangan kelima yang disisipkan adalah suku ke-6: $5+5\cdot6=35$.
\item[\textbf{5.}] \textbf{(E)} $a=18$, $r=\dis\frac{6}{18}=\dis\frac13$, $|r|<1$. $S_\infty=\dis\frac{18}{1-\frac13}=\dis\frac{18}{\frac23}=27$.
\item[\textbf{6.}] \textbf{(A)} $b=\dis\frac{37-17}{9-4}=4$ dan $a=U_4-3b=5$. $S_{12}=\dis\frac{12}{2}\big(2(5)+11(4)\big)=6(54)=324$.
\item[\textbf{7.}] \textbf{(E)} $r^{3}=\dis\frac{U_5}{U_2}=\dis\frac{96}{12}=8\Rightarrow r=2$, sehingga $a=\dis\frac{U_2}{r}=6$. $S_7=6\cdot\dis\frac{2^{7}-1}{2-1}=6(127)=762$.
\item[\textbf{8.}] \textbf{(C)} Aritmetika dengan $a=4.000.000$, $b=250.000$, $n=24$. $S_{24}=\dis\frac{24}{2}\big(2(4.000.000)+23(250.000)\big)=12(13.750.000)=165.000.000$.
\item[\textbf{9.}] \textbf{(D)} $\dis\frac{12}{1-r}=16\Rightarrow1-r=\dis\frac34\Rightarrow r=\dis\frac14$. $U_4=12\left(\dis\frac14\right)^{3}=\dis\frac{12}{64}=\dis\frac{3}{16}$.
\item[\textbf{10.}] \textbf{(E)} $U_{12}=S_{12}-S_{11}=(2\cdot144+60)-(2\cdot121+55)=348-297=51$. (Rumus umum: $U_n=4n+3$.)
\item[\textbf{11.}] \textbf{(C)} Dari $S_5=55$: $a+2b=11$. Dari $S_{10}=185$: $2a+9b=37$. Diperoleh $b=3$, $a=5$. Maka $U_{11}=35$ dan $U_{15}=47$, sehingga jumlahnya $\dis\frac52(35+47)=205$. (Cara cepat: jumlah blok lima suku membentuk barisan aritmetika $55,130,205$ dengan beda $25b=75$.)
\item[\textbf{12.}] \textbf{(B)} $b=\dis\frac{11-(-5)}{4}=4$, $a=-9$, sehingga $U_n=4n-13$. (1) $U_3=-1$ dan $U_4=3$: \textbf{benar}. (2) $S_{10}=5(-18+36)=90\neq96$: \textbf{salah}. (3) $4n-13=2.027\Rightarrow n=510$ (bulat): \textbf{benar}. (4) Selisih dua suku selalu kelipatan $4$, sedangkan $30$ bukan kelipatan $4$: \textbf{salah}.
\item[\textbf{13.}] \textbf{(C)} Rata-rata $n$ suku pertama $=\dis\frac{S_n}{n}=a+\dis\frac{(n-1)b}{2}$. Maka $a+4{,}5b=20$ dan $a+9{,}5b=30$, sehingga $b=2$, $a=11$. $S_{30}=\dis\frac{30}{2}\big(2(11)+29(2)\big)=15(80)=1.200$.
\item[\textbf{14.}] \textbf{(B)} $\dis\frac{a}{1-r}=36$ dan $\dis\frac{a}{1-r^{2}}=27$. Dibagi: $1+r=\dis\frac{36}{27}=\dis\frac43\Rightarrow r=\dis\frac13$, lalu $a=36\left(\dis\frac23\right)=24$. $U_3=ar^{2}=\dis\frac{24}{9}=\dis\frac83$.
\item[\textbf{15.}] \textbf{(D)} $U_2=a+b$, $U_5=a+4b$, $U_{14}=a+13b$. Syarat geometri: $(a+4b)^{2}=(a+b)(a+13b)\Rightarrow3b^{2}=6ab\Rightarrow b=2a=8$ (karena $b\neq0$). $S_{10}=\dis\frac{10}{2}\big(2(4)+9(8)\big)=5(80)=400$. (Cek: $12,36,108$ berrasio $3$.)
\item[\textbf{16.}] \textbf{(C)} Lapisan ke-$k$ berisi $\dis\frac{k(k+1)}{2}$ buah. Total $=\dis\frac12\left(\sum k^{2}+\sum k\right)=\dis\frac{n(n+1)(n+2)}{6}$. Untuk $n=15$: $\dis\frac{15\cdot16\cdot17}{6}=680$.
\item[\textbf{17.}] \textbf{(E)} Selisih antarsuku: $2,4,6,8,\ldots$ (aritmetika). Maka $U_n=3+\sum_{k=1}^{n-1}2k=3+n(n-1)$. $U_{20}=3+20\cdot19=383$.
\item[\textbf{18.}] \textbf{(D)} Misal ada $2k+1$ suku dan suku tengah $U_t$. Suku ganjil berjumlah $(k+1)U_t$, suku genap berjumlah $kU_t$. Maka $\dis\frac{k+1}{k}=\dis\frac{56}{48}=\dis\frac76\Rightarrow k=6$. Banyak suku $=2(6)+1=13$ (suku tengah $U_t=8$).
\item[\textbf{19.}] \textbf{(A)} Rasionalkan: $\dis\frac{1}{\sqrt k+\sqrt{k+1}}=\sqrt{k+1}-\sqrt{k}$. Jumlahnya bersifat teleskopik: $(\sqrt2-\sqrt1)+(\sqrt3-\sqrt2)+\cdots+(\sqrt{100}-\sqrt{99})=\sqrt{100}-1=9$.
\item[\textbf{20.}] \textbf{(D)} Sisi persegi berikutnya $=\dis\frac{s}{\sqrt2}$, jadi rasio sisi (dan keliling) adalah $\dis\frac{1}{\sqrt2}$. Keliling pertama $=32$. $S_\infty=\dis\frac{32}{1-\frac1{\sqrt2}}=\dis\frac{32\sqrt2}{\sqrt2-1}=32\sqrt2(\sqrt2+1)=64+32\sqrt2$. (Bila yang dijumlahkan luasnya: $\dis\frac{64}{1-\frac12}=128$.)
\end{enumerate}
\par\Needspace{6\baselineskip}\vspace{0.6em}{\bfseries\color{navy}\normalsize Pembahasan Isian Singkat}\par\vspace{0.2em}
\small\begin{enumerate}[leftmargin=2em,itemsep=2pt,topsep=2pt]
\item[\textbf{1.}] \textbf{Jawaban: $381$.} Ada $5+1=6$ interval: $r^{6}=\dis\frac{192}{3}=64\Rightarrow r=2$. Barisan: $3,6,12,24,48,96,192$ (7 suku), $S_7=3(2^{7}-1)=381$.
\item[\textbf{2.}] \textbf{Jawaban: $568$.} Suku-sukunya $13,16,19,\ldots,58$ ada $20-5+1=16$ suku. $S=\dis\frac{16}{2}(13+58)=8(71)=568$. (Cara lain: $\sum_{1}^{20}-\sum_{1}^{4}=590-22$.)
\item[\textbf{3.}] \textbf{Jawaban: $16.200$.} $8$ jam $=4$ periode: $200\cdot3^{4}=200\cdot81=16.200$.
\item[\textbf{4.}] \textbf{Jawaban: $17.157$.} Bilangannya $105,112,\ldots,497$. Banyak suku $n=\dis\frac{497-105}{7}+1=57$. $S_{57}=\dis\frac{57}{2}(105+497)=57\cdot301=17.157$.
\item[\textbf{5.}] \textbf{Jawaban: $\dis\frac{14}{11}$.} $1+0{,}27+0{,}0027+\cdots=1+\dis\frac{0{,}27}{1-0{,}01}=1+\dis\frac{27}{99}=1+\dis\frac{3}{11}=\dis\frac{14}{11}$.
\end{enumerate}
\par\Needspace{8\baselineskip}\vspace{0.6em}{\bfseries\color{navy}\normalsize Pembahasan dan Pedoman Penskoran Uraian}\par\vspace{0.2em}
\small\begin{enumerate}[leftmargin=2em,itemsep=3pt,topsep=2pt]
\item[\textbf{1.}] (a) $a+3b=14$ dan $S_8=4(2a+7b)=124\Rightarrow2a+7b=31$. Substitusi $a=14-3b$: $28-6b+7b=31\Rightarrow b=3$, $a=5$. (b) $U_n=5+(n-1)3=3n+2=83\Rightarrow n=27$. (c) $S_n=\dis\frac n2\big(10+3(n-1)\big)=\dis\frac{n(3n+7)}{2}$. $S_{24}=12\cdot79=948<1.000$ dan $S_{25}=\dis\frac{25\cdot82}{2}=1.025>1.000$. Jadi $n$ terkecil adalah $25$. \\ \textit{Penskoran: (a) 3 poin; (b) 2 poin; (c) 3 poin.}
\item[\textbf{2.}] (a) $S_4=S_2+r^{2}S_2=S_2(1+r^{2})\Rightarrow120=12(1+r^{2})\Rightarrow r^{2}=9\Rightarrow r=3$ (suku positif). $S_2=a(1+r)=4a=12\Rightarrow a=3$. (b) $U_n=3\cdot3^{n-1}=3^{n}$, sehingga $U_6=729$. (c) $S_n=\dis\frac{3(3^{n}-1)}{2}>1.000\Rightarrow3^{n}>\dis\frac{2.000}{3}+1\approx667{,}7$. Karena $3^{5}=243$ dan $3^{6}=729$, maka $n=6$ (cek: $S_5=363$, $S_6=1.092$). \\ \textit{Penskoran: (a) 3 poin; (b) 2 poin; (c) 3 poin.}
\item[\textbf{3.}] Tinggi pantulan ke-$n$: $h_n=10\left(\dis\frac35\right)^{n}$. (a) $h_3=10\cdot\dis\frac{27}{125}=2{,}16$ m. (b) Lintasan $=10+2(h_1+h_2+\cdots)$ dengan $h_1=6$, $r=\dis\frac35$: $\sum h_n=\dis\frac{6}{1-\frac35}=15$. Jadi total $=10+2(15)=40$ m. (c) $h_4=10\cdot\dis\frac{81}{625}=1{,}296$ m $>1$ dan $h_5=10\cdot\dis\frac{243}{3125}=0{,}7776$ m $<1$. Jadi pada pantulan ke-$5$. \\ \textit{Penskoran: (a) 2 poin; (b) 4 poin; (c) 2 poin.}
\item[\textbf{4.}] (a) Untuk $n\ge2$: $U_n=S_n-S_{n-1}=3n^{2}-n-\big(3(n-1)^{2}-(n-1)\big)=6n-4$. Untuk $n=1$: $U_1=S_1=2$, sama dengan $6(1)-4$, sehingga rumus berlaku untuk semua $n$. (b) $U_n-U_{n-1}=(6n-4)-(6n-10)=6$ (konstan), sehingga aritmetika dengan $a=2$ dan $b=6$. (c) $U_{11}+\cdots+U_{20}=S_{20}-S_{10}=(1.200-20)-(300-10)=1.180-290=890$. \\ \textit{Penskoran: (a) 3 poin; (b) 2 poin; (c) 3 poin.}
\item[\textbf{5.}] Tuliskan $\dis\frac xr,\ x,\ xr$. (a) Hasil kali $x^{3}=64\Rightarrow x=4$. (b) Setelah suku kedua ditambah $1$: $\dis\frac4r,\,5,\,4r$ aritmetika, sehingga $2(5)=\dis\frac4r+4r\Rightarrow2r^{2}-5r+2=0\Rightarrow(2r-1)(r-2)=0$. Karena barisan naik, $r=2$ (bukan $\dis\frac12$). (c) Barisan geometri: $2,4,8$. Barisan aritmetika: $2,5,8,\ldots$ dengan $a=2$, $b=3$. $S_{10}=\dis\frac{10}{2}\big(2(2)+9(3)\big)=5(31)=155$. \\ \textit{Penskoran: (a) 2 poin; (b) 3 poin; (c) 3 poin.}
\end{enumerate}
\par\vspace{0.6em}\begin{tcolorbox}[colback=softbg,colframe=navy!30,boxrule=0.5pt,arc=2pt,left=6pt,right=6pt,top=3pt,bottom=3pt]\footnotesize\textbf{Catatan penyusunan.} Soal disusun mengacu pada tipe-tipe yang lazim muncul di UTBK-SNBT/SBMPTN (pernyataan benar-salah, barisan aritmetika yang berubah menjadi geometri, deret tak hingga berparameter, jumlah suku ganjil-genap, deret teleskopik, dan isian singkat). Seluruh angka dan konteks dibuat orisinal, bukan salinan soal asli.\end{tcolorbox}
\end{document}
